\section{Analytical results}

The paper focuses on the six configurations in Table
\ref{tab:axm-regime-map}. This keeps the role of each elasticity separate.
Research complementarity, \(\sigma_{HM}<1\), is left for an extension.

\begin{table}[H]
\centering
\small
\setlength{\tabcolsep}{4pt}
\caption{Analytical map}
\label{tab:axm-regime-map}
\begin{tabular}{P{0.14\textwidth}P{0.235\textwidth}P{0.235\textwidth}P{0.235\textwidth}}
\toprule
Research & \(\sigma_{XL}<1\) & \(\sigma_{XL}=1\) & \(\sigma_{XL}>1\)\\
\midrule
\(\sigma_{HM}=1\) &
Labor bottleneck on a regular unbounded-\(A\) limit &
Candidate BGP governed by \(D(\omega_M)\) &
No finite-rate path with unbounded \(A\); the singular form is not characterized\\
\(\sigma_{HM}>1\) &
Same production bottleneck; \(s_M\to1\) if \(Aw\to\infty\) &
Automated candidate BGP governed by \(D(1)\) &
No finite-rate path with unbounded \(A\); conditional singular ratios are characterized\\
\bottomrule
\end{tabular}
\end{table}

\subsection{The price of research automation}

The research block can be understood before solving the dynamics. Dividing the
two conditional demands in \eqref{eq:axm-research-demands} gives
\begin{equation}
 \frac{AM}{H}
 =\left(\frac{\omega_M}{\omega_H}\right)^{\sigma_{HM}}
  (Aw)^{\sigma_{HM}},
 \label{eq:axm-input-ratio}
\end{equation}
and therefore
\begin{equation}
 \frac{s_M}{1-s_M}
 =\left(\frac{\omega_M}{\omega_H}\right)^{\sigma_{HM}}
  (Aw)^{\sigma_{HM}-1}.
 \label{eq:axm-automation-odds}
\end{equation}
The relevant relative price is \(Aw\): the wage relative to the normalized cost
of one unit of automated research services.

\begin{proposition}[Research automation]
\label{prop:axm-research-automation}
Along any interior path on which \(Aw\to\infty\), the automated contribution
to research satisfies
\[
 s_M\to
 \begin{cases}
  0,&\sigma_{HM}<1,\\
  \omega_M,&\sigma_{HM}=1,\\
  1,&\sigma_{HM}>1.
 \end{cases}
\]
For \(\sigma_{HM}>1\), automated research services asymptotically dominate the
research CES even though human research can remain positive at every finite date.
\end{proposition}

The proposition concerns the service \(AM\), not raw compute \(M\). This
distinction is important. A rising \(AM/H\) can reflect more research compute,
better capability, or both.

\subsection{Complementarity in final production}

When \(\sigma_{XL}<1\), the integrated monopolist does not supply an unlimited
ratio \(X/L\) as marginal cost tends to zero. On the regular interior branch, the
zero-marginal-cost limit of \eqref{eq:axm-monopoly-foc} sets marginal revenue to
zero. Hence
\begin{equation}
 s_X\longrightarrow
 \bar s_X=\frac{1-\sigma_{XL}}{1-\alpha\sigma_{XL}}\in(0,1).
 \label{eq:axm-complement-share}
\end{equation}
A constant CES share implies a finite constant \(X/L\).

\begin{proposition}[Labor bottleneck]
\label{prop:axm-complements}
Suppose \(\sigma_{XL}<1\), \(A\to\infty\), the monopoly allocation remains on the
regular interior branch,
\[
 \frac{1}{A(K/L)^\alpha}\to0,
\]
and the economy converges to a finite-rate balanced path. Then \(X/L\) converges
to a finite constant and final production becomes
asymptotically proportional to \(K^\alpha L^{1-\alpha}\). Consequently,
\[
 g_Y=g_K=g_L=n,\qquad g_{Y/N}=0,\qquad g_w=0.
\]
\end{proposition}

This is a conditional long-run result, not an existence theorem. In particular,
the declining marginal value of capability may lead the developer to a research
corner before \(A\) becomes unbounded. What the proposition establishes is that
unbounded AI capability does not by itself sustain per-capita growth when labor
and AI services are complements.

\subsection{The Cobb--Douglas final-production benchmark}

Set \(\sigma_{XL}=1\), and define
\begin{equation}
 \beta=(1-\alpha)\omega_X,\qquad
 \lambda=(1-\alpha)\omega_L,\qquad
 \nu=\frac{\beta}{\lambda}=\frac{\omega_X}{\omega_L}.
 \label{eq:axm-cd-definitions}
\end{equation}
The monopoly condition implies
\begin{equation}
 \frac{U}{Y}=\beta^2,
 \label{eq:axm-inference-share-cd}
\end{equation}
and the reduced production equation is
\begin{equation}
 Y^{1-\beta}
 =K^\alpha L^\lambda(\beta^2A)^\beta.
 \label{eq:axm-reduced-cd}
\end{equation}
If \(L\) and \(K,Y\) grow at their balanced rates, this equation gives
\begin{equation}
 g_Y=n+\nu g_A.
 \label{eq:axm-output-growth-cd}
\end{equation}

It is useful to treat Cobb--Douglas research and asymptotically automated research
separately. Let
\begin{equation}
 \bar s=
 \begin{cases}
  \omega_M,&\sigma_{HM}=1,\\
  1,&\sigma_{HM}>1\text{ and }Aw\to\infty.
 \end{cases}
 \label{eq:axm-sbar}
\end{equation}
Define the feedback denominator
\begin{equation}
 D(\bar s)=1-\eta\bar s(1+\nu).
 \label{eq:axm-feedback-denominator}
\end{equation}

\begin{proposition}[Balanced-growth candidates at unit elasticity]
\label{prop:axm-cd-bgp}
Suppose \(\sigma_{XL}=1\), research remains interior, its expenditure share has a
positive limit, and \(L/N\) converges to a positive constant.

\begin{enumerate}[label=(\roman*)]
\item If \(\sigma_{HM}=1\), every balanced-growth equilibrium must have
 \(D(\omega_M)>0\).
\item If \(\sigma_{HM}>1\), \(Aw\to\infty\), and the path converges to an
 asymptotically automated balanced-growth limit, that limit must have \(D(1)>0\).
\end{enumerate}

Whenever the relevant denominator is positive, the candidate growth rates are
\begin{equation}
 g_A=\frac{\eta n}{D(\bar s)},\qquad
 g_Y=g_K=g_C=n+\nu g_A,\qquad
 g_w=\nu g_A.
 \label{eq:axm-cd-growth-rates}
\end{equation}
Let \(\gamma=\nu g_A\). Its candidate output shares are
\begin{align}
 u&\equiv\frac{U}{Y}=\beta^2,\\
 i&\equiv\frac{\dot K+\delta K}{Y}
 =\frac{\alpha(n+\gamma+\delta)}{\rho+\delta+\gamma},\\
 m&\equiv\frac{M}{Y}
 =\frac{\beta^2\eta\bar s\,g_A}
 {\rho-n+(1-\eta\bar s)g_A},\\
 c&\equiv\frac{C}{Y}=1-u-i-m.
 \label{eq:axm-cd-shares}
\end{align}
A feasible interior candidate additionally requires
\[
 \rho>n,\qquad
 \rho-n+(1-\eta\bar s)g_A>0,\qquad c>0,
\]
the monopoly second-order condition, a feasible labor allocation, and the
transversality conditions.
\end{proposition}

The proposition states necessary growth restrictions and lists the remaining
conditions needed for existence; it does not infer existence from \(D>0\) alone.
The intuition for \(D\) is direct. A one-percent increase in \(A\) raises
production and hence research compute through \(\nu\). It also raises the services
produced by every unit of research compute directly. The term \(1+\nu\) combines
these two channels. Only the automated contribution \(\bar s\) transmits them
through the research CES, and \(\eta\) scales their effect on capability
improvement.

Because \(D(1)<D(\omega_M)\), gross substitution in research tightens the
balanced-growth restriction. With the benchmark values used below,
\(D(\omega_M)=0.8031\) and \(D(1)=0.4375\); both are positive. Removing an
additional \(A^\phi\) term is decisive for this conclusion: the baseline does not
add a separate knowledge-spillover channel to the feedback already carried by
\(AM\).

\subsection{Gross substitution in final production}

Let
\begin{equation}
 \theta=\frac{1-\alpha}{\alpha}.
 \label{eq:axm-theta}
\end{equation}
When \(\sigma_{XL}>1\), AI services dominate the final CES as their relative
quantity rises. The monopoly condition then implies
\begin{equation}
 \frac{U}{Y}\to(1-\alpha)^2,\qquad
 \frac{Y}{K}\sim \mathcal D A^\theta
 \label{eq:axm-high-sigma-scaling}
\end{equation}
for a positive constant \(\mathcal D\).

\begin{proposition}[No finite-rate steady state with gross substitution]
\label{prop:axm-no-bgp}
Suppose \(\sigma_{XL}>1\), \(A\to\infty\), and the equilibrium remains on an
interior monopoly branch on which \(s_X\to1\). Then there is no steady state with
finite constant growth rates. In particular, \(Y/K\to\infty\), so the gross return
to capital and household consumption growth become unbounded.
\end{proposition}

The proposition is a nonexistence result. It does not, by itself, prove that every
initial condition reaches a singularity: research can shut down, a constraint can
bind, or a different monopoly branch can intervene. The next result characterizes
the accelerating branch when it exists.

\begin{conditionalresult}[AI-dominated singular limit]
\label{prop:axm-singular-limit}
Assume \(\sigma_{XL}>1\), \(\sigma_{HM}>1\), \(0<\eta<1\), and an
interior equilibrium approaches a finite date at which \(A\to\infty\),
\(s_X\to1\), and \(s_M\to1\). Suppose its consumption, investment, and
research-compute shares, \(g_A/(Y/K)\), and \(qA/K\) converge to positive finite
limits, and \(L/N\) converges to a positive limit. Assume these convergences are
asymptotically regular: the growth rate of each convergent ratio, divided by
\(Y/K\), tends to zero. Write \(z=Y/K\) and
\[
 \Delta=1+\theta-\eta,\qquad
 h=\frac{\eta\alpha}{\Delta}.
\]
Then
\begin{align}
 \frac{g_A}{z}&\to h,&
 \frac{U}{Y}&\to(1-\alpha)^2,\\
 \frac{\dot K+\delta K}{Y}&\to\alpha-\theta h,&
 \frac{M}{Y}&\to\frac{\eta(1-\alpha)^2}{\Delta},\\
 \frac{qA}{K}&\to\frac{(1-\alpha)^2}{\eta\alpha},&
 \frac{C}{Y}&\to
 1-(1-\alpha)^2-\alpha+\theta h
 -\frac{\eta(1-\alpha)^2}{\Delta}.
 \label{eq:axm-singular-ratios}
\end{align}
Moreover,
\begin{equation}
 \dot z\sim\theta h z^2.
 \label{eq:axm-z-law}
\end{equation}
Thus the limiting differential equation reaches \(z=\infty\) in finite time.
Labor's production share tends to zero, while
\begin{equation}
 \frac{g_w}{z}\to
 \frac{\alpha-(\sigma_{XL}-1)h}{\sigma_{XL}}>0
 \label{eq:axm-wage-limit}
\end{equation}
which is positive if and only if
\(\sigma_{XL}<1/(\alpha\eta)\).
\end{conditionalresult}

The result characterizes a postulated asymptotic equilibrium branch; it does not
establish that the branch is reached from arbitrary initial conditions. Within
that branch, the model delivers a sharp separation when
\(1<\sigma_{XL}<1/(\alpha\eta)\): labor's share vanishes, but the wage eventually
rises. Above the threshold, the same conditional scaling implies an eventually
falling wage. This final case should not be interpreted without separately
checking monopoly optimality and convergence to the postulated branch.
